\documentclass[12pt]{article}
\usepackage{amsmath, amssymb, geometry, setspace}
\geometry{margin=1in}
\setstretch{1.2}
\setlength{\parindent}{0pt}



\title{Quiz 4\\ \vspace{10pt}
\large ECON 320 -- Introduction to Mathematical Economics}
\author{Instructor: Muhebullah Karimzada}
\date{October 22, 2025}
\begin{document}

\maketitle



\section*{Part A – Multiple Choice Questions (1 point each)}

\begin{enumerate}
\item The derivative of \( \ln(x^2 + 1) \) is  
\begin{enumerate}
\item \( \dfrac{2x}{x^2 + 1} \)
\item \( 2x\ln(x) \)
\item \( \dfrac{1}{x^2 + 1} \)
\item \( 2(x^2 + 1) \)
\end{enumerate}

\item The derivative of \( e^{2x} \) is  
\begin{enumerate}
\item \( e^{2x} \)
\item \( 2e^{2x} \)
\item \( 2xe^{2x} \)
\item \( e^{x^2} \)
\end{enumerate}

\item If \( y = e^{3x-5} \), then \( \frac{dy}{dx} = \)
\begin{enumerate}
\item \( e^{3x-5} \)
\item \( 3e^{3x-5} \)
\item \( 5e^{3x-5} \)
\item \( e^{3x} - 5 \)
\end{enumerate}
\newpage

\item The natural logarithm function \( \ln(x) \) is defined for  
\begin{enumerate}
\item all real \(x\)
\item \(x>0\)
\item \(x<0\)
\item \(x\neq0\)
\end{enumerate}

\item If \( f(x) = \ln(e^{2x}) \), then \( f'(x) = \)
\begin{enumerate}
\item \( 2e^{2x} \)
\item \( 2x \)
\item \( 2 \)
\item \( e^{2x} \)
\end{enumerate}

\item The elasticity of \( y = a e^{bx} \) with respect to \(x\) equals  
\begin{enumerate}
\item \( bx \)
\item \( b \)
\item \( e^{bx} \)
\item \( \frac{b}{x} \)
\end{enumerate}

\item The function \( y = \ln(x) \) is  
\begin{enumerate}
\item increasing and concave  
\item decreasing and convex  
\item increasing and convex  
\item decreasing and concave
\end{enumerate}
\end{enumerate}

\newpage

\section*{Part B – Analytical Problems (13 points total)}

\textbf{Question 1 (4 points)}\\
The utility of income \(x\) is \( U(x) = \ln(x + 2) \).

\begin{enumerate}
\item Compute \(U'(x)\) and \(U''(x)\).  
\item Show that utility is increasing but concave.  
\item Interpret the economic meaning of concavity in this context.
\end{enumerate}

\vspace{1em}

\textbf{Question 2 (4 points)}\\
A firm’s production function is given by \( Q(L) = 20L^{0.5}e^{-0.05L} \).

\begin{enumerate}
\item Derive the first-order condition for maximizing output.  
\item Find the optimal labor input \(L^*\).  
\item Verify that it corresponds to a maximum.  
\item Interpret the exponential term economically.
\end{enumerate}

\vspace{1em}

\textbf{Question 3 (5 points)}\\
The demand for a product is \( Q(P) = 100e^{-0.2P} \).

\begin{enumerate}
\item Write the revenue function \( R(P) = PQ(P) \).  
\item Find the price \(P^*\) that maximizes revenue.  
\item Compute the corresponding maximum revenue \(R(P^*)\).  
\item Verify the second-order condition.  
\item Explain the trade-off between price and quantity in exponential demand.
\end{enumerate}


\end{document}
